\documentclass[10pt,a4paper]{amsart}
\usepackage[margin=19mm]{geometry}
\usepackage{amsmath,amssymb,mathtools}
\usepackage[scr=rsfs,cal=euler]{mathalfa}
\usepackage{xfrac}
\usepackage[unicode,hidelinks,bookmarks=false]{hyperref}
\newcommand{\ie}{i.e.\ }
\newcommand{\cf}{cf.\ }
\newcommand{\resp}{respectively\ }
\newcommand{\Subsection}{\subsection}
\providecommand{\nobreakdash}{\nobreak\hbox{-}}
\setlength{\emergencystretch}{2em}
\multlinegap0pt
\usepackage{amssymb}
\usepackage{enumerate}
\usepackage[normalem]{ulem}
\usepackage{cancel}
\usepackage{tikz}
\usepackage{mathtools}
\newtheorem{lemma}{Lemma}
\newcommand{\R}{\mathbb{R}}
\title{Shape optimization of metastable states}
\author{No\'e Blassel, Tony Leli\`evre, Gabriel Stoltz}
\date{Source: arXiv:2507.12575v3 (CC BY 4.0)}
\begin{document}
\maketitle

\noindent\emph{Source and licence.} Shape optimization of metastable states. Version arXiv:2507.12575v3 (CC BY 4.0), distributed by arXiv under the Creative Commons Attribution 4.0 International licence, http://creativecommons.org/licenses/by/4.0/, as recorded in the arXiv licence metadata for this exact version and shown on the abstract page https://arxiv.org/abs/2507.12575v3. Full licence notice: \texttt{licenses/arxiv-2507.12575-CC-BY-4.0.txt}.\par\smallskip

\noindent\textit{Editorial scope.} Counted unit: the article's lemma lemma:unique\_maximizer (source0.tex lines 784--790). The article's complete original proof of it is delivered with it (source0.tex lines 791--810, one \texttt{proof} environment of 20 lines). No mathematics is written, repaired or re-derived: the statement and the proof are the article's own lines, in the article's own notation, and no retained line is substituted or re-flowed.  No further source-local context is needed: every cross-reference of every retained line resolves inside the two delivered spans.  Nothing was omitted from any retained span, so the package carries no omission record.  No retained line contains a citation, so the wrapper carries no bibliography and no bibliography entry.  No retained line contains a figure environment, an image-inclusion command or drawing code, so no image is delivered and the package declares no asset dependency.  Editorial formatting only, in addition: an AMS \texttt{amsart} preamble that restates the article's own declarations and macros used by the retained lines, namely the environments \texttt{lemma} on separate counters, together with the standard packages those lines need.  The retained environments print in this document as Lemma 1 in order of appearance, whereas the article prints the same items with the numbering of its own full document.  The complete native source of the article as distributed in the arXiv e-print for this exact version is bundled unchanged as \texttt{sources/181711/source0.tex}.
\begin{lemma}
    \label{lemma:unique_maximizer}
    Let~$B$ be the closed unit ball in a finite-dimensional Hilbert space $E$, and let $f:B\to\R$ be a concave function such that~$\underset{\theta\in B}{\sup}\, f(\theta)>0$, and is furthermore positively homogeneous of degree~$\alpha>0$. Then, there exists a unique maximizer $\theta^*\in \partial B$ for the problem
    \begin{equation}
        \underset{\theta\in B}{\sup}\,f(\theta).
    \end{equation}
\end{lemma}

\begin{proof}
    Since $B$ is compact, there exists a maximizer. Assume for the sake of contradiction the existence of two distinct maximizers~$\theta_1\neq \theta_2$.
    We let~$\theta = \frac12(\theta_1+\theta_2)$, and note that~$\|\theta\|_E<1$.

    First,
    \begin{equation}
        \begin{aligned}
            f\left(\theta\right)&\geq \frac12\left(f(\theta_1)+f(\theta_2)\right)\\
            &=\max_B\,f>0,
        \end{aligned}
    \end{equation}
    since~$f$ is concave, then
    \begin{equation}
        \begin{aligned}
        f(\theta/\|\theta\|_{E})&=\|\theta\|^{-\alpha}_{E}f(\theta)>f(\theta)\\
        &\geq \max_B\,f
        \end{aligned}
    \end{equation}
    using homogeneity. We have reached a contradiction, therefore there exists a unique maximizer $\theta^*$, which necessarily satisfies~$\|\theta^*\|=1$, since~$f(\theta/\|\theta\|_E)>f(\theta)$ whenever~$f(\theta)>0$, using homogeneity once again.
\end{proof}

\end{document}
